% --- Archon physics formalization source begin ---
% archon:physics
% archon:covers PhyXMiniProblems/problem_phyx_mini_0093.lean
% archon:source-report reports/phyx_mini/problem_phyx_mini_0093.source.json

\chapter{Physics problem phyx\_mini\_0093}
\label{ch:PhyXMiniProblems_problem_phyx_mini_0093}

\paragraph{Problem source.}
\#\# Physical scenario

A beam of unpolarized light of intensity \$I\_0\$ passes through a series of ideal polarizing filters with their polarizing axes turned to various angles as shown in figure.

\#\# Question

If we remove the middle filter, what will be the light intensity at point C?

\#\# Answer choices

- A: A: I = 0

- B: B: I = I0

- C: C: I = I0/2

- D: D: I=I0 cos²(90°)

\#\# Figure caption (auxiliary; use the image as primary evidence)

The image consists of a sequence of three polarizing filter diagrams, labeled A, B, and C, from left to right. 

1. **Scene A:**
   - An unpolarized light beam labeled \textbackslash{}( I\_0 \textbackslash{}) is shown entering a vertical polarizing filter with its axis vertically aligned.
   - The light is marked as "Unpolarized" before hitting the filter.
   - There are no additional angles or labels within the polarizer.

2. **Scene B:**
   - The second polarizer is also vertically aligned, with an angle labeled \textbackslash{}( 60\textasciicircum{}\textbackslash{}circ \textbackslash{}) shown in blue, indicating the angle of the polarizing axis relative to some reference.
   - This scene is labeled with a dot and the letter \textbackslash{}( B \textbackslash{}).

3. **Scene C:**
   - The third polarizer is similar but is labeled with an angle of \textbackslash{}( 90\textasciicircum{}\textbackslash{}circ \textbackslash{}).
   - Again, a dot and the letter \textbackslash{}( C \textbackslash{}) provide a label for this scene.

Each polarizer in the diagrams appears as an oval shape with a line through the center representing the polarizing axis. The progression suggests a sequential passing of light through each polarizer, each with a different orientation angle.

\#\# Dataset metadata

Wave Optics; Physical Model Grounding Reasoning, Implicit Condition Reasoning

\paragraph{Recorded answer/context.}
A: A: I = 0

\paragraph{Figure/image path.}
/root/proposal\_for\_physic/science-mango/phyx\_mini\_run/phyx\_data/test\_image/93.png

\paragraph{Formalization target.}
create a compiling Lean file with sorry bodies at `PhyXMiniProblems/problem\_phyx\_mini\_0093.lean`. The Lean declarations must preserve the physical quantities, dimensions or dimensional roles, figure labels, governing-law hypotheses, and final relation expressed by this problem.
Use Mathlib/Physlib names found through LeanExplore where available. If a physics API is missing, introduce faithful local abstractions rather than scalar placeholder aliases.

\begin{theorem}[Physics formalization target]
\label{thm:physics:phyx_mini_0093:target}
\lean{PhyXMiniProblems.ProblemPhyXMini0093.intensity_at_C_after_removing_middle_filter}
\uses{def:physics:phyx-mini-0093:phyxminiproblems-problemphyxmini0093-threepolarizersetup, def:physics:phyx-mini-0093:phyxminiproblems-problemphyxmini0093-threepolarizerfigurereadouts, def:physics:phyx-mini-0093:phyxminiproblems-problemphyxmini0093-threepolarizeropticslaws}
In the modified arrangement with the middle \(60^\circ\) polarizer removed,
unpolarized light passes first through the vertical polarizer A and then
through the \(90^\circ\) polarizer C.  Their crossed axes make the intensity
at C zero, answer A.
\end{theorem}
\begin{proof}
The ideal-polarizer law at A makes its outgoing beam linearly polarized along
axis angle \(0\), with half the incident intensity.  In the modified path the
next transition is directly through C, whose axis angle is \(\pi/2\).
Malus's law therefore multiplies the A intensity by
\(\cos^2(\pi/2-0)=0\).  Thus the scalar readout of the dimensioned output
intensity is zero.  Extensionality for the dimension-tagged quantity, together
with its common-unit representation, identifies the physical intensity itself
with zero.
\end{proof}

% --- Archon declaration topology begin ---
\section{Declaration topology}
\begin{definition}[optical Intensity Dimension]
  \label{def:physics:phyx-mini-0093:phyxminiproblems-problemphyxmini0093-opticalintensitydimension}
  \lean{PhyXMiniProblems.ProblemPhyXMini0093.opticalIntensityDimension}
  The physical dimension of optical intensity (power per unit area), namely M T⁻³ in the mass--length--time basis
\end{definition}
\begin{proof}
  This is the declared model definition of optical Intensity Dimension.
\end{proof}
\begin{definition}[Optical Intensity]
  \label{def:physics:phyx-mini-0093:phyxminiproblems-problemphyxmini0093-opticalintensity}
  \lean{PhyXMiniProblems.ProblemPhyXMini0093.OpticalIntensity}
  \uses{def:physics:phyx-mini-0093:phyxminiproblems-problemphyxmini0093-opticalintensitydimension}
  Optical intensity measured in one fixed common unit. The Physlib WithDim tag prevents this physical quantity from being identified with a bare real number, while .val is its scalar readout in that common unit
\end{definition}
\begin{proof}
  This is the declared model definition of Optical Intensity.
\end{proof}
\begin{definition}[Polarization State]
  \label{def:physics:phyx-mini-0093:phyxminiproblems-problemphyxmini0093-polarizationstate}
  \lean{PhyXMiniProblems.ProblemPhyXMini0093.PolarizationState}
  The polarization information carried by a beam in this problem. A linear polarization axis is represented by its real-valued radian angle from the vertical reference line in the figure
\end{definition}
\begin{proof}
  This is the declared model definition of Polarization State.
\end{proof}
\begin{definition}[Light Beam]
  \label{def:physics:phyx-mini-0093:phyxminiproblems-problemphyxmini0093-lightbeam}
  \lean{PhyXMiniProblems.ProblemPhyXMini0093.LightBeam}
  \uses{def:physics:phyx-mini-0093:phyxminiproblems-problemphyxmini0093-opticalintensity, def:physics:phyx-mini-0093:phyxminiproblems-problemphyxmini0093-polarizationstate}
  A light beam together with its measured optical-intensity readout and polarization state. All intensity readouts use one fixed common physical unit, so the real scalar records a measurement rather than identifying intensity as an untyped scalar primitive
\end{definition}
\begin{proof}
  This is the declared model definition of Light Beam.
\end{proof}
\begin{definition}[Ideal Linear Polarizer]
  \label{def:physics:phyx-mini-0093:phyxminiproblems-problemphyxmini0093-ideallinearpolarizer}
  \lean{PhyXMiniProblems.ProblemPhyXMini0093.IdealLinearPolarizer}
  An ideal linear polarizing filter and the orientation of its transmission axis
\end{definition}
\begin{proof}
  This is the declared model definition of Ideal Linear Polarizer.
\end{proof}
\begin{definition}[Passes Through Ideal Linear Polarizer]
  \label{def:physics:phyx-mini-0093:phyxminiproblems-problemphyxmini0093-passesthroughideallinearpolarizer}
  \lean{PhyXMiniProblems.ProblemPhyXMini0093.PassesThroughIdealLinearPolarizer}
  \uses{def:physics:phyx-mini-0093:phyxminiproblems-problemphyxmini0093-lightbeam, def:physics:phyx-mini-0093:phyxminiproblems-problemphyxmini0093-ideallinearpolarizer}
  The ideal-filter transition law. Unpolarized incident light loses half its intensity. Linearly polarized incident light obeys Malus's cosine-squared law. In either case, the outgoing polarization is aligned with the filter axis
\end{definition}
\begin{proof}
  This is the declared model definition of Passes Through Ideal Linear Polarizer.
\end{proof}
\begin{definition}[Three Polarizer Setup]
  \label{def:physics:phyx-mini-0093:phyxminiproblems-problemphyxmini0093-threepolarizersetup}
  \lean{PhyXMiniProblems.ProblemPhyXMini0093.ThreePolarizerSetup}
  \uses{def:physics:phyx-mini-0093:phyxminiproblems-problemphyxmini0093-opticalintensity, def:physics:phyx-mini-0093:phyxminiproblems-problemphyxmini0093-lightbeam, def:physics:phyx-mini-0093:phyxminiproblems-problemphyxmini0093-ideallinearpolarizer}
  Physical objects and labeled beam locations in the three-polarizer figure. beamAtCWithoutB belongs to the modified experiment in which the middle filter is removed; its intensity is deliberately not fixed by this structure
\end{definition}
\begin{proof}
  This is the declared model definition of Three Polarizer Setup.
\end{proof}
\begin{definition}[Three Polarizer Figure Readouts]
  \label{def:physics:phyx-mini-0093:phyxminiproblems-problemphyxmini0093-threepolarizerfigurereadouts}
  \lean{PhyXMiniProblems.ProblemPhyXMini0093.ThreePolarizerFigureReadouts}
  \uses{def:physics:phyx-mini-0093:phyxminiproblems-problemphyxmini0093-threepolarizersetup}
  Direct readouts from the source label and the displayed filter-axis geometry. The three angles are respectively 0°, 60°, and 90° from vertical
\end{definition}
\begin{proof}
  This is the declared model definition of Three Polarizer Figure Readouts.
\end{proof}
\begin{definition}[Three Polarizer Optics Laws]
  \label{def:physics:phyx-mini-0093:phyxminiproblems-problemphyxmini0093-threepolarizeropticslaws}
  \lean{PhyXMiniProblems.ProblemPhyXMini0093.ThreePolarizerOpticsLaws}
  \uses{def:physics:phyx-mini-0093:phyxminiproblems-problemphyxmini0093-passesthroughideallinearpolarizer, def:physics:phyx-mini-0093:phyxminiproblems-problemphyxmini0093-threepolarizersetup}
  Applications of the ideal-polarizer law to the original A--B--C arrangement and to the modified arrangement with B removed. These transition assumptions do not state the requested intensity at C
\end{definition}
\begin{proof}
  This is the declared model definition of Three Polarizer Optics Laws.
\end{proof}
% --- Archon declaration topology end ---
% --- Archon physics formalization source end ---
